\chapter{语言任务的求解与执行}
\label{chap:nlp-solving-and-execution}

“这次读书会改到什么时候？”需要依据通知回答；“把我的预约改到通知中的时间”却要求另一种交付物：预约记录真的发生变化，而且没有改错对象。第\ref{chap:nlp-question-answering-dialogue}章已经讨论如何取得证据、澄清要求和维护交流状态。本章继续追问：要求明确以后，怎样得到可检查的解答，怎样交付符合规格的程序，又怎样确认环境中的任务已经完成？

这三个问题共用候选形成、执行反馈和修订，但验收依据不同。算术题要核对表达式、单位和原始条件，程序任务要核对约定输入上的行为，环境任务则要读取行动后的实际状态。把这些依据分开，才能定位错误究竟来自题意映射、计算、程序实现，还是执行期间变化的环境。以下数值、程序仓库、工具接口和运行轨迹均为教学构造；真实论文只提供方法依据，不为这些构造案例的效果背书。

\section{数学与逻辑问题求解}
\label{sec:nlp-execution-solving}

静态求解通常可以先固定题目，再反复构造和检查候选，不必改变外部业务状态。不过，“算一个金额”“找一组满足排班规则的取值”和“证明一个命题”的输出并不相同。我们先建立所求对象，再比较候选的形成方式，最后选择与对象相符的核验器。

\subsection{条件、目标与形式化对象}
\label{sec:nlp-execution-formalization}

语言中的量、对象和条件需要成为计算能够读取的输入。第\ref{chap:nlp-text-analysis}章讨论的指代、语义关系和时间规范化在这里提供映射依据，但本章还要记录每个形式对象来自哪一项条件。缺少单价属于信息不足，把“不打折”读成“打折”属于解析错误，而在已正确建立的约束上没有找到解才属于求解问题；三者不能混成一个“答错”标签。

\subsubsection{表达式树与带类型的变量}
\label{sec:nlp-execution-expression}

考虑一道独立于读书会通知的数量题：“购买3本书，每本25元；另买5支笔，每支4元。书打八折，笔不打折，不收其他费用。应付多少元？”一个可检查的表示应同时保存数值、单位和修饰范围。设书的单价为$p_{\mathrm{b}}=25$元/本、数量为$n_{\mathrm{b}}=3$本，笔的单价为$p_{\mathrm{p}}=4$元/支、数量为$n_{\mathrm{p}}=5$支，书的支付比例为$r_{\mathrm{b}}=0.8$。比例没有单位，“八折”指支付原价的80\%，不是减去原价的80\%。

\term{表达式树}（expression tree）把已知量作为叶，把运算作为内部节点；根节点给出待求金额。这里先分别计算书与笔的金额，只对书的分支乘折扣，再相加：
\begin{equation}
  a=(p_{\mathrm{b}}n_{\mathrm{b}})r_{\mathrm{b}}
      +p_{\mathrm{p}}n_{\mathrm{p}}
    =(25\times3)\times0.8+4\times5=80\text{元}.
  \label{eq:nlp-execution-price}
\end{equation}
这棵树可写为$\operatorname{Add}(\operatorname{Mul}(\operatorname{Mul}(25,3),0.8),\operatorname{Mul}(4,5))$；每个数字另带原文位置和单位标签。构造器按“单价乘数量得到金额”“金额才能相加”建立节点，既可以由规则完成，也可以用题目与目标树配对监督条件生成模型；推断时从候选树恢复括号和运算次序。

类型检查能排除$25\text{元/本}+3\text{本}$，却排除不了$(25\times3+4\times5)\times0.8=76$元，因为后一式的单位完全正确，只是把折扣错误地施加给了笔。若将“另买5支”错绑定为“买3支”，还会得到72元。因此，恢复树以后要逐项核对数量指代、折扣作用域和“无其他费用”这项封闭条件，不能只检查能否代入数字。

\subsubsection{有限域约束与回溯求解}
\label{sec:nlp-execution-csp}

排班题不能总用一个算式直接算出答案。\term{约束满足问题}（constraint satisfaction problem, CSP）为每个变量指定有限候选域，再要求这些变量的联合取值满足全部条件。构造一个两时段值班例子：甲、乙、丙各选时段1或2；甲乙不能同时值班，乙必须早于丙，丙只能选时段2。令变量$a,b,c$分别表示三人的时段，便得到
\[
 D_a=D_b=\{1,2\},\quad D_c=\{2\},\qquad a\ne b,\quad b<c.
\]
本题求所有可行赋值，没有额外的最小成本目标，也没有待训练参数。若使用布尔变量$z_i$表示人员是否被选择，“至少一人”对应$\sum_i z_i\geq1$，“恰好一人”对应$\sum_i z_i=1$，“甲乙不能同时被选”对应$z_{\mathrm{a}}+z_{\mathrm{b}}\leq1$。这些语句限定的是不同集合，不能互换。

\term{前向检查}（forward checking）是在每次赋值以后，删除尚未赋值变量中与当前赋值冲突的取值。固定变量顺序$a,b,c$，域内按升序尝试，并为每个分支保存一份域快照。令$a=1$后，$D_b$只剩$\{2\}$；再令$b=2$，约束$b<c$会把$D_c=\{2\}$删空。搜索回到$a$处，恢复原来的候选域，改试$a=2$，此时$D_b=\{1\}$，随后$b=1,c=2$满足全部条件。包括失败尝试在内，共进行了$a=1,b=2,a=2,b=1,c=2$五次赋值。

继续枚举可确认唯一解为$(2,1,2)$；去掉$a\ne b$后，$(1,1,2)$也会成为解，而增加$a<b$则无解。若只找到一个解便停止，只能报告“找到可行排班”，不能据此声称唯一。前向检查不在所有未赋值变量之间反复传播；更强的一致性算法可能提前发现$b=2$不可能，本例的五次操作不是CSP求解的固有成本。

24点可以用类似的状态搜索，但状态不是人员取值。给定数字多重集$\{3,3,8,8\}$，每个数都须使用且只能使用一次，允许二元加、减、乘、除和括号。状态保存“剩余值、产生该值的表达式、已消耗的原始位置集合”。每次选两个不重叠的表达式，合为一个值；减法与除法要保留两个顺序，零不能作除数。用有理数保存$8/3$，可以避免浮点误差把最终的24判成近似失败。仅保存剩余数值可用于某些可达性去重，但交付表达式仍需保留父指针或原始数字的使用记录。

\subsubsection{模理论约束与Z3接口}
\label{sec:nlp-execution-z3}

若约束涉及整数算术与布尔组合，可以将公式交给\term{可满足性模理论}（satisfiability modulo theories, SMT）求解器：它判断是否存在某种解释，同时满足逻辑连接词和所用数域的规则。Z3是这种接口的一个实现。下面用官方指南的SMT-LIB命令表达同一个排班题，声明类型、加入断言、检查可满足性，再读取模型。

\begin{codeblock}[label={code:nlp-execution-z3}]{排班约束的SMT-LIB表示}
(declare-const a Int)
(declare-const b Int)
(declare-const c Int)
(assert (and (<= 1 a) (<= a 2)))
(assert (and (<= 1 b) (<= b 2)))
(assert (= c 2))
(assert (distinct a b))
(assert (< b c))
(check-sat)
(get-model)
\end{codeblock}

这里\code{Int}限定整数，而不是允许1.5这样的中间时段。按前面的枚举，公式应为可满足，模型只能给出$a=2,b=1,c=2$。调用器只有在收到\code{sat}后才读取模型；\code{unsat}表示提交的公式无模型，\code{unknown}表示这次检查未能作出决定，不能归为无解。可用\code{push}/\code{pop}临时增加条件、检查后恢复原断言集。这是求解器的作用域接口，不等于上一小节的手写前向检查算法。相关命令依据见本章文献说明中的Z3指南。

代码\ref{code:nlp-execution-z3}用于说明约束接口，实际使用时须固定Z3版本并保存完整输出，将返回模型与独立枚举或原约束交叉核验。即使返回\code{sat}，也应把模型重新代回原题：遗漏\code{distinct}会允许甲乙同时值班，求解器没有收到这条限制，便无从替作者补上。

\subsection{候选形成与搜索}
\label{sec:nlp-execution-candidates}

形式化对象确定以后，还要决定把多少计算投入哪些候选。直接生成一个答案、逐步写解答、独立采样多条路径以及搜索共享前缀，分别改变了候选空间的组织方式。比较时应固定题目、模型、提示示例和预算，记录生成数量、有效候选比例、评分调用与最终选择结果。

\subsubsection{步骤条件生成：CoT}
\label{sec:nlp-execution-cot}

直接输出“80元”几乎没有留下定位错误的中间材料。\term{思维链}（chain of thought, CoT）提示在题目和答案之间加入可读步骤，使后续生成可以利用前面的数量关系与计算。提示示例提供“题目、分步解答、答案”的格式，模型参数沿用预训练或已有微调结果；只改变提示不等于重新训练模型。原始CoT工作的提示方法见文献说明。

令$x$为题目和示例，$z_1,\ldots,z_m$为步骤片段，$y$为答案。生成分布可按前缀展开为
\begin{equation}
 p_{\theta}(z_1,\ldots,z_m,y\mid x)
 =\left[\prod_{i=1}^{m}p_{\theta}(z_i\mid x,z_1,\ldots,z_{i-1})\right]
 p_{\theta}(y\mid x,z_1,\ldots,z_m).
 \label{eq:nlp-execution-cot}
\end{equation}
式中的片段内部仍按词元自回归生成，$\theta$为固定模型参数；步骤边界由提示格式或后处理确定。解码使用贪心选择或预先约定的采样策略，到答案结束标记或长度上限停止，而不是执行一个自动成立的数学推导规则。

数量题的一条正确路径是“书原价75元；书折后60元；笔20元；合计80元”。若第三步写成“笔也打折，计16元”，就能把错误定位到“不打折”的条件，而不必猜测最终76元的来历。修订时重新提供该条件与错误步骤，从受影响位置重写并复算后续金额，不能只把末尾76改成80。可读步骤是可检查的外部产物，不是形式证明，也不足以断言它忠实暴露了模型内部计算。

\subsubsection{多路径答案聚合：自一致性}
\label{sec:nlp-execution-self-consistency}

不同步骤顺序可以到达同一个金额，例如先算每本折后20元，或先算三本原价75元再打折。\term{自一致性}（self-consistency）利用这种汇合：从固定模型采样多条完整路径，提取并规范化最终答案，再选择出现次数最多的答案。原文的路径采样与答案边际聚合见文献说明；本节将单位规范化和空结果处理明确为教学接口。

固定预算为六条，构造输出为“80元”“80.0元”“人民币80元”“76元”“76.00元”和“无法确定”。将前三条映射为$(80,\text{元})$，中间两条映射为$(76,\text{元})$，末条标为无有效答案；有效率为$5/6$，票数为3比2。若答案用“分”，应先按明确换算规则转成元；缺单位且无法从问题恢复时，不盲目与金额合并。全无有效答案时返回未求得答案，平票时采用预先规定的独立核验或弃权策略，不事后挑选有利的平局规则。

设规范化函数为$N$，有效路径索引集合为$V$，投票结果是
\[
 \hat y=\argmax_y\sum_{i\in V}\mathbb{I}[N(y_i)=y].
\]
这里每条路径算一票，即使两条文本接近也不自动减少权重；相关错误因此可能形成多数。例如六条路径中四条都把笔打折，就会稳定选出76元。验证器重排的best-of-$N$则在同一候选集中按可接受性分数选最高者，不按出现次数选择。两种方案要分别报告最终正确率和生成、规范化、评分的总成本，六次尝试的结果不能作为单次生成结果呈现。

\subsubsection{状态扩展与回退：Tree of Thoughts}
\label{sec:nlp-execution-tot}

独立生成完整路径会重复相同前缀，也难在中途更换一个局部决定。\term{思维树}（Tree of Thoughts, ToT）把若干步骤组成候选状态，以扩展器提出下一步，以评价器估计哪些状态值得继续，再由搜索控制器分配预算。扩展器和评价器都可以通过提示调用固定语言模型，不要求另行训练；评价可输出分值，也可比较候选后投票。原文分别使用宽度受限的广度优先搜索和可回退的深度优先搜索，具体方法见文献说明。

为了单独检查搜索控制，下面固定一棵教学候选树，不把手设评分冒充模型测量。根$R=\{3,3,8,8\}$提出两个分支：$A=\{3,3,16\}$来自$8+8$，评分0.9；$B=\{3,8,8/3\}$来自$8/3$，评分0.8。$A$再提出$C=\{6,16\}$（0.9）与$D=\{3,19\}$（0.4）；$B$提出$E=\{8,1/3\}$（0.8），其中$1/3=3-8/3$。已有表达式随状态保存。两数状态枚举加、乘、两个方向的减法与除法，共六个终态，用精确规则检查是否为24。

\begin{algorithm}[宽度受限的ToT广度优先搜索]
\label{alg:nlp-execution-tot-bfs}
\begin{algorithmic}[1]
\Require 初态$s_0$，扩展器$G$，评分器$V$，宽度$b$，深度$T$，预算$B$
\Ensure 一个已核验解，或附搜索范围的未找到状态
\State $F\gets\{s_0\}$；初始化预算计数与访问记录
\For{$t=1,\ldots,T$}
  \State 对$F$中状态依次调用$G$；预算不足则返回未找到
  \State 汇集合法子状态$C$；按足够表达约束的状态键去重
  \State 为$C$逐项评分，记录评分成本
  \If{$C$中存在通过完整规则核验的终态}
    \State \Return 按预定平局规则选择的解及其表达式
  \EndIf
  \State 删除错误终态；按分数降序、生成次序升序排序
  \State $F\gets$排序后至多$b$个非终态
  \If{$F=\varnothing$}
    \State \Return 当前候选范围内未找到
  \EndIf
\EndFor
\State \Return 深度上限内未找到
\end{algorithmic}
\end{algorithm}

令$b=2,T=3$。第一层保留$A,B$；第二层比较$C,D,E$，保留$C,E$，$D$因宽度限制被剪掉。第三层对$C$的六种运算都无法得到24，但$E$有
\[
 \frac{8}{3-8/3}=\frac{8}{1/3}=24.
\]
这里分母不是$3-8$再除以3，恢复括号是解码的一部分。以“一次展开一个非终态”为一次扩展，本轨迹展开$R,A,B,C,E$，共5次；生成和评分条目数为$2+3+12=17$。保留宽度最多为2，评分前的候选池仍可有12项，所以不能把宽度2误报为只存两个对象。

\begin{algorithm}[带回溯的ToT深度优先搜索]
\label{alg:nlp-execution-tot-dfs}
\begin{algorithmic}[1]
\Require 当前状态$s$、深度$t$、扩展器$G$、评分器$V$、阈值$\tau$与共享预算
\Ensure 已核验解或本分支未找到
\If{预算耗尽或$t=T$}
  \State \Return 未找到
\EndIf
\State $C\gets G(s)$；删除非法项，按状态键去重并计入预算
\State 对$C$评分；按分数降序、生成次序升序排序
\For{排序后的$c\in C$}
  \If{$c$通过完整终态核验}
    \State \Return $c$及其表达式
  \ElsIf{$c$为非终态且$V(c)\geq\tau$且未被当前路径访问}
    \State $r\gets\operatorname{DFS}(c,t+1)$
    \If{$r$为解}
      \State \Return $r$
    \EndIf
  \EndIf
\EndFor
\State \Return 未找到；调用者恢复父状态，继续下一个候选
\end{algorithmic}
\end{algorithm}

以根深度0、$\tau=0.1$运行相同候选树，DFS先走$R,A,C$，$C$无解便回退到$A$，再试$D$；$D$也无解，才回到根走$B,E$。它展开6个非终态，生成和评分$2+2+6+6+1+6=23$项。本实现先评分一个节点的全部子项，再寻找终态，因此找到解的那一批六项也全部计费。把逐项检查改成遇到解立即停止，会得到另一组成本，比较时必须说明。

两个结果只反映这棵固定候选树与手设排序。真实ToT的生成可能漏掉$B$，评价可能错误剪掉$E$；此时“未找到”不是原题无解。可达性相同的数值状态能够共享缓存，但若问题还限制操作次数、表达式形式或可用数字身份，去重键必须包含这些信息。这里的BFS在同层只保留前$b$项，控制方式与束搜索相近；A*或蒙特卡洛树搜索可以作为其他预算分配方案，不能写成原始ToT实验已经采用的算法。

\subsubsection{程序辅助求解：PAL}
\label{sec:nlp-execution-pal}

算式关系已经写对时，仍可能在75乘0.8这样的计算上出错。\term{程序辅助语言模型}（Program-aided Language Models, PAL）将中间关系写成程序，把数值运算交给解释器。提示示例由题目和带有语义变量名的程序组成，解释可以放在注释中；生成程序以后执行，再从指定变量提取答案。PAL原文使用的主要流程是生成完毕后的单次执行，并非默认每行执行后再让模型续写。

代码\ref{code:nlp-execution-pal}是按数量题构造、实际在本地Python解释器执行的教学候选，没有调用外部语言模型。它将折扣显式绑定到书的金额上，使用十进制定点语义避免二进制浮点显示误差。

\begin{codeblock}[language=Python,label={code:nlp-execution-pal}]{金额题的程序候选}
from decimal import Decimal

book_price = Decimal("25")
book_count = 3
book_pay_ratio = Decimal("0.8")
pen_price = Decimal("4")
pen_count = 5
book_payment = book_price * book_count * book_pay_ratio
pen_payment = pen_price * pen_count
answer = book_payment + pen_payment
print(f"{answer:.2f}")
\end{codeblock}

实际标准输出为\code{80.00}。结果提取器检查程序成功结束、输出只有一个合法金额，再附上原题单位形成“应付80元”。若变量名拼错，异常属于执行失败；若金额变量绑定到了错误数量，程序可能无异常却回答错误；若使用错误折扣，解释器仍会忠实输出76。需要保存程序、解释器版本、退出码与输出，才能区分这些情形。真实生成代码还应置于限制权限、时间和资源的隔离环境；本章只执行已核读的本地教学程序。

\subsection{结果验证与解释}
\label{sec:nlp-execution-verification}

候选形成后，需要分别检查最终答案、中间过程和解释的对应关系。一条解答可能碰巧给出80元，却包含错误等式；也可能程序算对了，随后的文字把折扣对象说错。核验时保留候选本身，错误定位回到具体条件或步骤，而不是另生成一段赞同候选的说理。

\subsubsection{计算回代与规则核验}
\label{sec:nlp-execution-rule-check}

规则核验器读取原题的变量表、允许运算与目标，再独立重算候选。金额题可用分为单位的整数关系交叉检查：书支付$2500\times3\times4/5=6000$分，笔支付$400\times5=2000$分，总计8000分。与PAL共用同一段错误公式只会重复错误；独立核验至少要从原题重新绑定折扣范围，并比较书、笔两项明细。

24点的检查分成语法、数值与资源使用三项。解析器只允许整数叶与四则二元运算，从树中收集叶值及出现次数，检查是否恰为$\{3,3,8,8\}$，再以有理数递归求值并拒绝除零。表达式$8\times3+3-3$的值确实是24，却用了三个3、一个8，应被拒绝；直接提交常数24也没有消耗题给的数字。搜索器保留的位置集合方便构造，独立核验器则从最终表达式重建证据，防止错误的状态记录被原样信任。

这种协议可用于预先明确规则的算术题，却不能照搬为开放写作的正确性判定。即使答案有可比的数字，也须先约定舍入、容差、单位及题目是否允许多个答案；否则“数值不同”可能只是表示不同，“数值相同”也可能掩盖对象不同。

\subsubsection{完整候选评分：ORM}
\label{sec:nlp-execution-orm}

自然语言解答不总能完全转成规则检查，但可以训练评分器帮助选择候选。\term{结果奖励模型}（outcome reward model, ORM）对题目和完整候选给出可接受性分数。GSM8K验证器的代表做法是先由生成器为训练题采样多条解答，按最终答案是否正确形成监督，再训练验证器；测试时生成多个候选，选择验证器评分最高者。来源见文献说明。

令训练样本为$(x_i,y_i,\ell_i)$，其中$\ell_i\in\{0,1\}$由答案核验得到，$v_{\phi}(x_i,y_i)$为评分器预测。一个结果监督目标为
\begin{equation}
 L(\phi)=-\sum_i\left[
 \ell_i\log v_{\phi}(x_i,y_i)
 +(1-\ell_i)\log(1-v_{\phi}(x_i,y_i))\right].
 \label{eq:nlp-execution-orm}
\end{equation}
这里学习的是评分参数$\phi$，不是本次候选中的算式；完整解答作为输入，最终位置的分数用于重排。实际实现可增加词元级辅助目标，本式只表达本节所需的结果二分类监督。

在另一组教学候选中，假设五条解答都给出76元，只有一条给出符合题意的80元。投票会选择76元，ORM却可能将正确候选排在首位；反过来，它也可能给措辞流畅的错误解答最高分。评价要分开统计候选集合是否包含正确解、正确候选能否排在错误候选之前，以及真正选出的那一条是否正确。训练、调参和测试按题目划分，并记录候选来自哪种生成器，避免验证器只记住训练题或熟悉的错误风格。只按答案生成的标签还会把“步骤错误但末尾碰巧写80”的候选标成正例。

\subsubsection{中间步骤评分：PRM}
\label{sec:nlp-execution-prm}

要定位这种内部错误，监督就需要进入步骤。\term{过程奖励模型}（process reward model, PRM）在每个步骤末，根据题目、已有步骤和当前步骤预测其质量。《Let's Verify Step by Step》以正、负、中性标签记录正确合理、错误不合理与存在歧义的步骤，再训练模型预测步骤标签；它与ORM的差别首先是监督对象，不是是否使用大语言模型。

令第$i$步标签为$\ell_i$，模型预测为$p_{\phi}(\ell_i\mid x,z_1,\ldots,z_i)$，训练累加所标注步骤的负对数似然。使用时可提取“正”类概率$q_i$，以
\begin{equation}
 v_{\mathrm{process}}(x,z)=\prod_{i=1}^{m}q_i
 \quad\text{或等价地}\quad
 \log v_{\mathrm{process}}=\sum_{i=1}^{m}\log q_i
 \label{eq:nlp-execution-prm}
\end{equation}
对完整候选排序。原文采用正标签概率的乘积。该汇总是排序规则；未经校准与条件假设，不能直接解释为整条证明成立的真实概率。步骤越多乘积通常越小，比较时还要关注切分方式；中性标签如何计分也须预先确定。

构造三步候选：“书原价75元；$75\times0.8=65$元；加20元，合计80元”。它的答案数值正确，第二步乘法和第三步加法却有错误。若过程评分为$0.99,0.10,0.98$，乘积为0.09702；另一条正确三步解答每步得0.9，乘积为0.729，排序便偏向后者。这是手设分数展示，不是PRM实测。可以用第一个低于阈值的步骤提示复核位置，但发现$75\times0.8=65$仍应交给算术检查；若把PRM放入搜索，每次前缀评分也必须计入调用成本。

\subsubsection{形式证明对象检查：Lean}
\label{sec:nlp-execution-lean}

步骤评分仍在预测判断。对于已经准确形式化的命题，可以改用证明检查：把命题表示为类型，把候选证明表示为该类型的一个项，检查器验证项的构造是否合法。Lean 4的这一接口见文献说明中的官方教材。本例只使用命题、函数应用和合取类型，不依赖额外数学库。

\begin{codeblock}[label={code:nlp-execution-lean}]{从已知前提构造证明项}
theorem eligibility
    (registered timely allowed : Prop)
    (rule : And registered timely -> allowed)
    (facts : And registered timely) : allowed :=
  rule facts

#print axioms eligibility
\end{codeblock}

将\code{registered}理解为已提交信息，\code{timely}理解为在截止时刻前提交，\code{rule}是明确给出的“同时满足两条件则允许”的前提；证明项把\code{facts}交给\code{rule}，得到\code{allowed}。模型若产生证明策略，策略也应最终生成可检查的证明项。验收时固定Lean 4具体版本、依赖和允许的公理，排除\code{sorry}、\code{admit}等未完成占位，并审查\code{\#print axioms}列出的依赖；不能仅凭没有语法错误接受文件。

代码\ref{code:nlp-execution-lean}用于说明命题与证明项的接口，实际使用时须记录\code{lean --version}并核验检查器输出。若把\code{timely}错误定义为“4月16日当天”，证明仍只会处理这个错误定义，无法恢复原通知严格的“4月16日18:00前”。形式化映射与证明对象检查因此是两项独立职责：前者对齐语言要求，后者确认在提交的定义和前提下结论成立。

\section{程序任务}
\label{sec:nlp-execution-programs}

PAL把程序当作求解题目的中间工具；程序任务则直接交付软件产物或对其行为的说明。用户可能要求解释现有函数、生成新函数、转换接口，或修复仓库中的错误。这些任务都需要固定代码版本、运行环境、依赖和允许修改的范围，才能知道解释针对哪份程序、测试实际执行了什么。

图\ref{fig:nlp-execution-program-loop}将开发反馈与独立验收分开。开发测试可以进入候选修订，独立验收只能在候选冻结后评价，否则它也成为开发信息。这里“隐藏”是相对于待评价系统而言，不要求教材读者永远看不到测试内容。

\begin{figure*}[htbp]
  \centering
  % Adapted from academic-drawing/templates/tikz_template.tex.
% Main ports: south -> north. Feedback remains left of the main column.
\begin{tikzpicture}[
  every node/.style={font=\figurefont\small,align=center,text=BookInk},
  block/.style={book node,font=\figurefont\small,text width=44mm,minimum height=11mm},
  flow/.style={book arrow,rounded corners=1.2mm},
  feedback/.style={flow,dashed,draw=BookAmber}
]
  \node[block,draw=FigureInput,fill=FigureInput!10] (spec)
    at (-45mm,0) {规格与允许变化\\输入、输出、保持条件};
  \node[block,draw=FigureInput,fill=FigureInput!10] (assets)
    at (45mm,0) {已有资产\\仓库版本、接口与依赖};
  \node[block,draw=FigureModel,fill=FigureModel!10] (candidate)
    at (0,-25mm) {生成候选或局部补丁\\保留版本与修改位置};
  \node[block] (test) at (0,-47mm)
    {在固定环境运行\\目标测试与相关回归};
  \node[block] (result) at (0,-69mm)
    {读取实际结果\\返回值、异常、失败断言};
  \node[block,text width=27mm,draw=BookAmber,fill=BookAmber!8]
    (revise) at (-61mm,-69mm) {定位并修订\\计入总预算};
  \node[block,draw=FigureOutput,fill=FigureOutput!10] (freeze)
    at (0,-91mm) {按预定规则选定\\冻结候选与环境};
  \node[block,text width=38mm] (stop) at (60mm,-69mm)
    {停止并报告\\预算耗尽或环境失败};
  \node[block,draw=FigureOutput,fill=FigureOutput!10] (heldout)
    at (0,-118mm) {独立验收\\留出测试与结果记录};
  \coordinate (inputs) at (0,-14mm);
  \draw[draw=BookMuted,line width=.7pt,rounded corners=1.2mm]
    (spec.south) |- (inputs);
  \draw[draw=BookMuted,line width=.7pt,rounded corners=1.2mm]
    (assets.south) |- (inputs);
  \draw[flow] (inputs) -- (candidate.north);
  \draw[flow] (candidate.south) -- (test.north);
  \draw[flow] (test.south) -- (result.north);
  \draw[feedback] (result.west) -- node[above] {未通过} (revise.east);
  \draw[feedback] (revise.north) |- (candidate.west);
  \draw[flow] (result.east) -- (stop.west);
  \draw[flow] (result.south) -- node[right] {通过} (freeze.north);
  \draw[flow] (freeze.south) -- (heldout.north);
  \draw[draw=BookMuted,line width=.8pt,dotted]
    (-78mm,-104mm) -- (80mm,-104mm);
  \node[anchor=east,align=right] at (79mm,-114mm)
    {不回流为\\开发反馈};
\end{tikzpicture}

  \caption{程序候选、实际执行与修订。实线给出从规格到验收的处理路径，虚线把开发失败反馈给修订；点线划开开发与独立评价。每个候选使用相同环境和测试协议，测试结论只覆盖约定的行为范围。}
  \label{fig:nlp-execution-program-loop}
\end{figure*}
\FloatBarrier

\subsection{程序理解与规格恢复}
\label{sec:nlp-execution-understanding}

“这个函数按通知发布时间排序吗？”不能只读函数名作答。程序理解需要取得与问题相关的代码，连接跨文件调用，再检查关键输入上的行为。由代码恢复的是当前行为的候选说明；判断这种行为是否符合原始意图，还需要需求、接口契约或有效测试等独立材料。

\subsubsection{符号检索与BM25、嵌入检索}
\label{sec:nlp-execution-code-retrieval}

对仓库提问时，检索单位宜包含路径、限定符号名、函数签名、注释和代码片段，而不只是任意字符窗口。固定索引的版本后，将问题“通知排序为何忽略同一天内的先后”转成符号查询\code{sort\_notices}以及“发布时间、日期解析、排序键”等词。名称匹配取得直接定义，BM25按词项检索候选，嵌入检索补充措辞不同但语义相关的辅助函数；排序机制回用第\ref{chap:nlp-classification-matching-organization}章。

假设仓库同时有\code{notices.py:sort\_notices}和测试文件中的同名辅助函数，仅按名称会召回两者。读取实际入口的导入关系，可以确认生产函数调用了\code{time\_utils.py:parse\_local}；再扩展到该定义，便能看到时间字段是否在解析中被截断。检索器若经监督训练，正例来自问题与相关符号的标注，负例应包括这种同名干扰项；若使用现成嵌入，参数来自预训练，仍须说明切片长度、候选数和上下文预算。

最终交给解释器或语言模型的材料应带路径、行号和版本定位，注明尚未读到的依赖。注释“按完整发布时间排序”只是意图线索；若辅助函数实际返回日期，解释就应报告同日记录只能保持输入次序。排名靠前的片段回答“可能在哪里”，跨文件证据才开始回答“为什么会这样”。

\subsubsection{AST与控制流、数据流证据}
\label{sec:nlp-execution-ast}

取得代码后，需要追踪条件和数据怎样影响返回值。\term{抽象语法树}（abstract syntax tree, AST）将源码解析为函数、条件、调用、运算与返回等语法节点。Python的\code{ast.parse}提供这类结构，官方接口见文献说明。对下列教学函数，AST显示一个空输入分支、一个类型检查分支，以及由\code{set}传给\code{sorted}的嵌套调用。

\begin{codeblock}[language=Python,label={code:nlp-execution-sort-trace}]{需要核对去重行为的排序函数}
def unique_sorted(values):
    if not values:
        return []
    if any(type(x) is not int for x in values):
        raise TypeError("integers required")
    return sorted(set(values))
\end{codeblock}

沿控制流看，空列表直接返回，不进入后续分支；非空输入会检查全部元素类型，不合格时抛出异常。沿数据流看，\code{values}先进入集合构造，重复值在排序之前已经被合并，所以注释若声称“保留重复项”，就与实际计算不符。类型判断用\code{type(x) is int}还意味着布尔值不被接受，尽管Python中的布尔类型属于整数的子类。

构造控制流图需要把条件的真假出口与返回、异常出口相连，数据流分析需要跟踪定义与使用；这些并非AST节点自动给出的全部语义。在跨文件例子中，若导入名被重新赋值或调用目标来自运行时注册表，静态调用边只能标为可能关系。AST也不能单独断定任意循环终止或两个程序等价；不确定边应进入解释的限制，而不是补写成确定事实。

\subsubsection{运行轨迹与边界输入}
\label{sec:nlp-execution-traces}

将具体输入绑定到程序后，可以记录实际返回、异常与关键中间变量。代码\ref{code:nlp-execution-sort-trace}对空列表返回\code{[]}，对\code{[3,1,3]}经集合\code{\{1,3\}}返回\code{[1,3]}，对\code{[1,"2"]}在类型检查处抛出\code{TypeError}。这三个输入分别覆盖空分支、重复项合并和异常分支，帮助把“排序并去重，只接受非空情况下的整数元素检查”这样的说明落到代码位置。

轨迹应区分运行时采集与模型预测。模型文字“第三行得到两个元素”可以引导核查，却不是程序真的执行到第三行的证据。教学脚本直接调用上述函数，捕获返回值与异常；对于更复杂的程序，可在允许的位置记录变量快照，或在调试器中观察相同输入，保留环境版本。

三个样本仍没有覆盖生成器输入、自定义容器或无穷序列；如果接口只允许整数列表，就应在规格中写清并据此设计测试。分支覆盖衡量走过哪些出口，不衡量输入空间是否穷尽。恢复的行为说明应保留这种输入范围，不能从空、重复和异常三个例子推出全语言对象上的保证。

\subsection{程序生成与转换}
\label{sec:nlp-execution-generation}

生成任务需要把规格变成候选程序，再决定接受哪一个。以通知排序为例，约定输入为记录列表，每条记录含唯一业务标识和带时区的发布时间；输出按时间升序，保留重复记录与同键的原次序，不修改输入；缺失时区、无效日期和不合格记录应失败而非猜测。规格同时决定正常、边界和无效输入，测试预期不从候选代码反推。

多个候选的评价不能与单次交付混用。HumanEval用函数签名、文档字符串和测试检查生成程序，pass@$k$回答“抽取$k$个候选时，至少一个通过测试的概率”。对一道题预先生成$n\geq k\geq1$个候选，其中$c$个通过测试。从这$n$个候选中不放回均匀选$k$个，共有$\binom nk$种子集，全部失败的子集有$\binom{n-c}{k}$种，所以该题的估计量为
\begin{equation}
 \widehat{\operatorname{pass@}k}
 =1-\frac{\binom{n-c}{k}}{\binom nk}.
 \label{eq:nlp-execution-pass-k}
\end{equation}
约定$n-c<k$时分子为0。若$c=0$则估计为0；$k=1$时退化为$c/n$。在固定采样分布下独立同分布生成候选，该对子集事件的平均是“$k$次采样至少一次成功”概率的无偏估计；跨题指标再对题目取平均。自适应修订得到的候选并不自动符合这一采样协议。

构造$n=5,c=2,k=2$，十个二元子集有三个全失败，因此pass@2为$7/10=0.7$。把$\hat p=2/5$代入$1-(1-\hat p)^2$却得到0.64，不能取代组合估计。更关键的是，若部署选择器固定挑中一个失败候选，本题实际交付仍然失败；0.7不是这个选择器单次成功率。应另报实际选定代码的测试结果、总采样数和执行预算。来源中的HumanEval协议见文献说明。

\subsubsection{条件代码生成与检索示例}
\label{sec:nlp-execution-conditional-code}

代码生成的输入可以组织为“任务描述、函数签名、类型、前后条件、允许依赖、示例和已有代码”。若用配对的需求与程序训练，目标是最大化参考代码在这些条件下的似然；直接使用已有代码模型时，参数来自预训练或既有微调，本次规格和示例只进入上下文。解码形成候选后依次解析、导入、调用并运行开发测试，解析失败的候选保留错误位置，不计为功能通过。

对通知排序，最短候选可能是\code{sorted(rows, key=...)}，真正的问题在排序键是否保留时刻。检索到相似的日志排序例子，可以复用“先规范化时区，再按时间排序”的接口用法；但需核对依赖版本和返回类型，不能把处理日期的示例照搬为处理时刻。若上下文太长，可先确定时间解析、输入验证与排序三个子任务，再实现相互一致的接口，而不是让每个子任务独立猜测数据类型。

开发断言至少包括空列表、同日逆序、相同时间的稳定次序、重复项计数以及输入未被修改。还可使用有独立依据的性质：排序前后记录多重集相同，输出时间单调不降，重复排序得到相同结果。仅检查“多重集不变”仍会放过完全没排序的程序，所以性质测试应互相补充；由候选自己生成且未经审查的断言，不能作为它自己的独立验收依据。

\subsubsection{语法约束解码：SynCode}
\label{sec:nlp-execution-syncode}

即使规格清楚，自由生成仍可能漏括号或在JSON字段之间遗漏逗号。SynCode把形式语言的语法接入解码，在每个前缀上计算哪些词元可以继续生成，再用掩码约束模型分布。这里上下文无关文法描述整体结构，终结符的正则语言由\term{确定有限自动机}（deterministic finite automaton, DFA）处理；语法和分词器来自接口定义，不是通过当前候选的测试学习出来的。

困难在于模型词元与语法终结符不一一对应。例如已经生成的前缀为
\begin{quote}
\code{\{"date":"2025-04-}
\end{quote}
一个词元可能解码为\code{18","}，同时补完日期字符串、产生逗号并开启下一个字段名；只看“下一个终结符是字符串”不够。SynCode增量解析已确定部分，保留尚未词法确定的后缀，枚举可接续的终结符序列，再结合后缀在DFA中的状态查询预计算的词元掩码。合并各条可接受路径的掩码后，将不允许词元的对数分数设为$-\infty$，归一化并采样或取最大值，随后更新解析状态。

保证必须按原文的方向理解。以$P$表示当前前缀，$A(P)$表示仍能扩成合法字符串的下一词元集合，$M(P)$表示掩码放行集合。原文称为soundness的性质对应$A(P)\subseteq M(P)$，即不错误排除合法续接；若枚举的接续终结符序列足够长，满足其定理关于长度大于词表内词元字符长度等条件，才得到反向包含。实现常用较短接续序列，可能保守放行非法续接。因此应保留完整解析、结束符和长度上限检查，不能把理论条件略去后声称所有输出都完整合法。

即使JSON完整解析，\code{"date":"2025-02-30"}仍是字符串而非有效日期；即使日期有效，也可能填的是发布日期而不是活动日期。结构约束减少的是特定语法错误，日期域、类型和任务对象仍由后续接口检查承担。这里的掩码机制依据SynCode正文，来源见本章文献说明。

\subsubsection{执行反馈修订：Self-Debugging}
\label{sec:nlp-execution-self-debugging}

候选能解析以后，测试才能揭示行为差异。Self-Debugging通过提示让模型阅读候选、解释代码，并结合执行结果或反馈修订代码；模型参数保持不变。原文区分有单元测试的反馈与主要依赖执行结果、解释的场景，其中某些“轨迹解释”由模型生成，并非运行时采集。本章使用真实测试返回作为外部反馈，再用解释组织定位。

设候选把“4月16日18:00前”实现为\code{submitted <= deadline}。在两个时间都保留时刻的条件下，输入恰为\code{2025-04-16T18:00:00+08:00}，实际返回\code{True}，而根据严格截止要求，预期为\code{False}。反馈记录应包含输入、预期、实际返回、相关代码和测试来源；“边界不对”过于含糊，“请改成小于”则已经直接给出补丁，不适合在比较自主修订能力时混入。

一次修订把比较符改为\code{<}，再运行相同失败用例和17:59允许、次日拒绝等回归用例。若日期解析器此前已丢失时刻，仅修改比较符会把17:59也拒绝；新反馈于是要求继续追踪调用链。下一小节组将完整展示这种跨文件情况。每轮保存候选差分、实际输出和累计调用数，达到预设修订上限仍失败时报告未解决。

无测试时，代码解释可以帮助发现与规格不符的字段或分支，但解释中的“现在正确”不增加独立证据。比较修订与重新采样时，应固定总生成词元、模型调用和测试预算，并报告最终选择规则，否则较多轮次带来的收益会被误归因于反馈机制本身。

\subsubsection{AST规则转换与条件代码转换}
\label{sec:nlp-execution-transformation}

转换任务的输入还包括原程序和允许变化：例如保持时刻排序行为，只把记录键\code{published\_at}改为新接口的\code{publishedAt}。AST规则可以只匹配目标函数中对这一字符串键的下标访问，修改相应节点并重新生成源码；全局文本替换则可能误改文档示例或另一个业务对象。转换后要检查语法、目标键是否替换完整，以及不在允许范围内的节点是否保持。

更一般的语言迁移可用“源代码、源语言语义、目标接口和示例”条件生成候选，再编译或解析目标程序，在同一组输入上比较返回、异常和副作用。规则转换的知识来自人工编写的匹配与重写条件；学习式转换的参数来自配对程序或代码预训练，检索示例只补充当前库接口。两者都要明确哪些行为必须保持，哪些变化是用户要求的。

跨语言比较不能只看运算符外观。例如对$-3$和2，向零截断的整数除法得到$-1$，Python的\code{-3 // 2}得到$-2$；用\code{/}则得到浮点数$-1.5$。固定宽度无符号32位加法可以将$2^{32}-1+1$回绕为0，Python整数则得到$2^{32}$，需要按目标契约决定是否显式取模。异常类型、时区库以及可变对象别名也可能改变行为。通过一组对照输入只支持该测试范围；代码文本接近原文，并不是语义保持的判据。

\subsection{程序修改与修复}
\label{sec:nlp-execution-repair}

仓库修复与从零生成的差别在于，大部分行为必须保持不变。SWE-bench以问题描述和既有仓库版本为输入，要求形成解决问题的补丁，强调了代码理解、跨文件依赖与执行验收的联系；本章不把补丁与参考文本相似当作修复成功。来源见文献说明。

本节构造两文件仓库，并明确补充业务接口：\code{eligible}只判断校外来宾提交信息是否早于截止时刻，校内人员直接返回允许；它不验证信息完整性，也不等同于现实中的完整报名系统。时间字符串必须带时区，比较前统一到北京时间。共同通知仍发布于2025年4月15日，活动\code{ML-0418}改至4月18日周五14:00、海棠楼201，校外提交截止为4月16日18:00；下面的程序不改变这些已知事实。

\subsubsection{分层定位与局部修复：Agentless}
\label{sec:nlp-execution-agentless}

Agentless将处理分为定位、修复、补丁验证三个职责。定位时先用目录结构筛选可疑文件，并以嵌入检索补充文件候选；再从类、函数和变量声明组成的文件骨架缩小到相关元素，最后读取这些元素的实际内容，确定编辑位置。修复阶段给出位置附近的上下文，生成局部搜索／替换补丁；流程由外部程序固定，不让模型自行决定下一种工具动作。这里遵循2024年修订版的方法正文，不把它简化成只有定位和修复的两阶段。

问题报告为“截止当刻仍被允许，同一天的通知排序也不正确”。从仓库目录定位到\code{notices.py}，看到\code{eligible}和\code{sort\_notices}；追踪两者使用的\code{parse\_local}，进一步定位到\code{time\_utils.py}。原程序的相关片段为：

\begin{codeblock}[language=Python,label={code:nlp-execution-repair-base}]{两个文件中相互关联的日期错误}
# time_utils.py
from datetime import datetime, timedelta, timezone
LOCAL = timezone(timedelta(hours=8))

def parse_local(value):
    timestamp = datetime.fromisoformat(value)
    if timestamp.tzinfo is None:
        raise ValueError("timezone required")
    return timestamp.astimezone(LOCAL).date()

# notices.py
from time_utils import parse_local

def eligible(kind, submitted_at, deadline):
    if kind == "internal":
        return True
    if kind != "external":
        raise ValueError("unknown participant kind")
    return parse_local(submitted_at) <= parse_local(deadline)

def sort_notices(rows):
    return sorted(rows,
                  key=lambda row: parse_local(row["published_at"]))
\end{codeblock}

这里有两条独立证据：\code{.date()}使09:00与14:00都变成4月15日，排序只能保留原输入次序；\code{<=}即使收到完整时刻，也会接纳恰好18:00的提交。正确修订需要去掉日期截断，并把报名比较改成\code{<}，排序函数自身无需修改。若定位器只提供\code{eligible}，后续修复器就看不到辅助函数的返回类型，容易产生下一小节中的回归补丁。

实际教学运行使用Python 3.9.6，每个候选在独立子进程导入自己的两文件副本，设置5秒执行超时，避免模块缓存把上一候选带入下一候选。候选及测试均已核读，只运行本地标准库代码。真实模型产生的补丁还需要限制文件范围、拒绝测试删改与无关环境修改；不能让候选通过改变验收条件来消除失败。

\subsubsection{多候选补丁与测试选择}
\label{sec:nlp-execution-patch-selection}

在相同基线生成三个教学候选：$P_1$只保留完整时刻，$P_2$只改严格比较，$P_3$同时修正两处。目标测试有两项：截止当刻应拒绝，同日通知应按时刻排序。七项回归分别检查截止前允许、次日拒绝、校内直接允许、空列表、重复项与稳定次序、缺时区报错以及未知人员类型报错。表\ref{tab:nlp-execution-patches}记录实际运行结果。

\begin{table}[htbp]
  \centering
  \begin{tabular}{lcc}
    \toprule
    候选 & 目标测试通过 & 回归通过 \\
    \midrule
    原程序 & 0/2 & 7/7 \\
    $P_1$：仅保留时刻 & 1/2 & 7/7 \\
    $P_2$：仅严格比较 & 1/2 & 6/7 \\
    $P_3$：两处共同修正 & 2/2 & 7/7 \\
    \bottomrule
  \end{tabular}
  \caption{跨文件日期修复的教学运行结果。分母固定，失败均为行为断言不符，不是导入或依赖故障。}
  \label{tab:nlp-execution-patches}
\end{table}
\FloatBarrier

$P_2$通过了截止当刻拒绝，却把17:59也变成与截止同一天的日期，导致严格小于为假。这说明“原失败测试转为通过”尚不足以选择补丁。教学选择器要求目标测试全部通过、约定回归全部保持，并检查差分只修改两处允许位置，因而选出$P_3$；若多个候选都通过，再按预定的修改范围与平局规则选择。将$P_1$收到的边界失败反馈接入修订，也会得到$P_3$，但这种顺序修订与独立多候选采样应作为不同协议记录。

原程序和三个候选共运行$4\times9=36$个开发用例。冻结$P_3$后另运行四个预先留出的验收用例：等价UTC偏移的截止时刻、截止前一微秒、无效日历日期和无效记录类型，四项全部通过。这是本地构造程序的功能核验，不是对Agentless或语言模型准确率的实验。测试文件、两文件候选、差分和JSON报告均由配套脚本保存，隐藏用例未反复用于修订。

Agentless原文还生成复现测试，先在原仓库上筛掉不能复现问题的测试，再结合原有回归测试选择补丁。其修订版先保留回归失败最少的候选，再使用选定复现测试过滤，并对规范化补丁作多数重排；复现测试不可靠、无人通过时有回退策略。本节要求全部指定回归通过的教学选择器更严格，不冒充原文算法。无论采用哪种规则，模型生成测试的预期都应对照问题独立审查，无法建立可靠验收依据时只报告候选与已验证范围。

\section{工具与环境任务}
\label{sec:nlp-execution-environment}

程序测试可以在固定副本上重跑，环境中的一次修改却可能影响后续所有操作。用户说“把我的预约改到通知中的时间和地点，不要发送通知，保留备注”，系统需要找到正确记录、读取当前状态、执行允许的修改，再核验结果。工具返回、用户补充与其他入口的修改都可能改变下一步，因此这里的基本单位是带观察的行动，而不是一次完整文本生成。

以下补充一组独立的教学业务状态，不声称共同通知含有这些资料。用户\code{U7}有个人预约\code{R17}，关联活动\code{ML-0418}，原时间为2025年4月17日14:00、个人预约地点字段为海棠楼101、备注为“带笔记”，版本为1。同一用户另有同名预约\code{R18}，关联其他活动，必须保持不变。共同通知给出的目标仍为4月18日周五14:00、海棠楼201；修改个人预约并不修改活动通知，也不办理报名。模拟环境用\code{Haitang-201}作为房间标识，所有时间带北京时间偏移。

图\ref{fig:nlp-execution-environment-loop}把语言目标、控制状态与实际环境分开。图中“观察”只收录工具真正返回的内容；模型预测“应该已经改好”仍是待验证判断。配套程序是运行在本地内存中的确定性教学模拟，没有调用真实预约系统，也没有使用语言模型生成动作。

\begin{figure*}[htbp]
  \centering
  % Adapted from academic-drawing/templates/tikz_template.tex.
% Main path uses south -> north ports; feedback is outside the main column.
\begin{tikzpicture}[
  every node/.style={font=\figurefont\small,align=center,text=BookInk},
  block/.style={book node,font=\figurefont\small,
    text width=43mm,minimum height=11mm},
  side/.style={block,text width=29mm},
  flow/.style={book arrow,rounded corners=1.2mm},
  recovery/.style={flow,dashed,draw=BookAmber}
]
  \node[block,draw=FigureInput,fill=FigureInput!10] (goal) at (0,0)
    {目标及保持条件\\对象、授权、目标版本};
  \node[side,draw=FigureInput,fill=FigureInput!10] (user) at (-61mm,0)
    {用户澄清\\或更正};
  \node[block,draw=FigureModel,fill=FigureModel!10] (plan) at (0,-24mm)
    {判断与计划\\依赖、预算、下一步};
  \node[block] (action) at (0,-48mm)
    {绑定参数并行动\\模式、身份、版本检查};
  \node[block,draw=FigureOutput,fill=FigureOutput!10] (environment)
    at (0,-72mm) {实际环境\\记录及副作用};
  \node[side] (external) at (61mm,-72mm)
    {其他入口\\修改共同状态};
  \node[block] (observation) at (0,-96mm)
    {真实工具观察\\字段、回执或异常};
  \node[block] (state) at (0,-120mm)
    {事件更新状态\\事实、假设、未知操作};
  \node[side,draw=BookAmber,fill=BookAmber!8] (recover) at (61mm,-120mm)
    {先回查旧请求\\恢复或重新读取};
  \node[block,draw=FigureOutput,fill=FigureOutput!10] (check) at (0,-144mm)
    {受控读取与结束核验\\目标及保持条件};
  \node[block,draw=FigureOutput,fill=FigureOutput!10] (done) at (0,-168mm)
    {输出完成证据\\或未完成与未知项};
  \draw[flow] (user.east) -- node[above] {更正} (goal.west);
  \draw[flow] (goal.south) -- (plan.north);
  \draw[flow] (plan.south) -- (action.north);
  \draw[flow] (action.south) -- (environment.north);
  \draw[flow] (external.west) -- (environment.east);
  \draw[flow] (environment.south) -- (observation.north);
  \draw[flow] (observation.south) -- (state.north);
  \draw[flow] (state.south) -- (check.north);
  \draw[flow] (check.south) -- (done.north);
  \draw[flow] (state.west) -- (-35mm,-120mm)
    -- (-35mm,-24mm) -- (plan.west);
  \node[text width=22mm,align=right,anchor=east] at (-39mm,-75mm)
    {更新信息\\重选行动};
  \draw[recovery] (state.east) --
    node[above,text width=17mm] {写入未知} (recover.west);
  \draw[recovery] (recover.north) -- (61mm,-96mm) -- (observation.east);
  \node[text width=29mm,align=center] at (61mm,-151mm)
    {回执描述已发生操作\\回读确认当前状态};
\end{tikzpicture}

  \caption{语言要求到环境结果的执行闭环。实线连接行动与真实观察；左侧反馈更新目标和重选行动，右侧虚线处理未知写入。用户更正影响目标，其他入口的修改影响环境，二者通过不同证据进入状态。}
  \label{fig:nlp-execution-environment-loop}
\end{figure*}
\FloatBarrier

\subsection{目标与工具接口}
\label{sec:nlp-execution-tools}

把任务整理为目标记录、所需字段、前置条件、保持条件和结束依据，可以避免“调用了更新接口”成为唯一成功标准。本例的前置条件是身份已确认、版本仍有效且目标时段可用；保持条件是备注和另一预约不变、不发送通知；结束依据是读取到\code{R17}的目标时间与地点，并检查保持条件。语言约束、接口说明与当前环境事实分别提供这些输入，不能由生成器随意补全。

\subsubsection{工具检索与动作选择}
\label{sec:nlp-execution-tool-selection}

工具目录可能同时有搜索预约、查询时段、读取记录、更新记录和发送通知。检索以“还缺什么信息或变化”为查询，而不是把整段用户要求直接匹配到最显眼的动词。当前尚无对象ID，应该先选搜索工具；已确认ID但不知道记录版本，则应选读取工具；目标与前置条件都已确认，更新工具才成为候选。发送通知虽与“预约修改”语义相关，却违反本例的否定约束。

可先按名称、说明和输入输出模式检索前$k$个候选，再由分类器在候选中选择，或让模型直接生成工具名。若训练分类器，监督来自任务状态与允许工具的配对；直接使用已有模型时，工具定义和少量示例进入上下文，参数不因这次调用而更新。相关性计算复用第\ref{chap:nlp-classification-matching-organization}章的方法，随后还须过滤权限、必需输入和副作用不相容的工具。工具名生成器输出目录外名称时，应报告不可用，而不是尝试拼出一个接口。

教学目录的\code{search}只读，支持按名称精确查找或按活动编号与用户查找；\code{read}读取单条记录，\code{availability}检查时段；\code{update}只修改时间和房间，不发送通知。它们能共同完成本例，但若真实更新接口强制给所有人发消息，就没有满足当前约束的动作，必须说明冲突或请求用户重新授权。把同义工具替换进计划之前，需要比较的正是这种契约，而不只是名称。

\subsubsection{函数参数生成与模式约束}
\label{sec:nlp-execution-tool-arguments}

动作确定以后，参数要从有来源的字段绑定，而不是只填满JSON。搜索结果提供\code{id}，刚读取的记录提供\code{expected\_version}，通知提供目标时间和房间，控制器生成任务内唯一请求号。一个正常更新请求为：

\begin{codeblock}[label={code:nlp-execution-tool-request}]{本地模拟工具的结构化参数}
{
  "id": "R17",
  "expected_version": 1,
  "start": "2025-04-18T14:00:00+08:00",
  "room": "Haitang-201",
  "request_id": "task7-goal1-v1"
}
\end{codeblock}

模式要求这五个字段齐全、不接受额外字段，版本为整数，其余字段为字符串。调用前先解析JSON、检查模式，再解析带时区时间、查询ID和版本，并把对象所属用户与活动编号同当前目标比对。自由文本解析可能把“版本一”转换错，模式约束能减少缺字段与类型错；第\ref{sec:nlp-execution-syncode}节的语法掩码只负责可生成的结构，不能替代对象绑定与日历验证。

本地接口核验中，\code{R404}得到不存在错误，字符串版本\code{"1"}得到模式错误，2月30日得到日期错误，过期版本得到冲突。另一组更隐蔽的候选把合法日期填为通知发布日期4月15日，或把ID填为同名的\code{R18}：这些值本身均可解析，若权限也允许，接口可以成功修改错误目标。因此调用器需在写前核对目标身份，结束判定器还要再次读取正确对象；参数文本和接口回执应分别保存。

\subsection{计划与行动选择}
\label{sec:nlp-execution-planning}

任务预算至少包括模型决策次数、工具调用次数与总时间；失败调用和回查也消耗预算。本例每次服务调用计一次，参数检查是本地计算，结束回读计入调用，独立评价器另行记账。计划必须为异常和结束核验留出空间，不能将预算全部花在修改动作上。一般状态转移与策略学习属于强化学习，本节只展开工具执行需要的控制接口。

\subsubsection{固定计划与计划—执行分工}
\label{sec:nlp-execution-plan-execute}

固定计划可写为“按名称找到预约，读取版本，查询目标时段，修改字段，回读核验”。规划器根据目标和工具说明给出子目标及依赖，执行器则把查询输出绑定到下一步参数；未取得ID时，后续计划中应保留“待搜索结果确定”，而不是预填\code{R17}。规划器可以是人工规则，也可以是使用任务示例的固定参数模型；计划质量以可执行性和目标达成为依据，不以步骤数量评分。

执行器对每一步应记录预期观察与停止条件。搜索预期恰好找到符合用户和活动约束的一条记录，多条则澄清，零条则修改查找方式；读取预期返回当前版本；时段查询预期为可用，否则重选经用户允许的候选；写入后必须读取实际字段。若执行器只返回“第五步失败”，规划器就难以区分时段已占用、版本冲突和记录不存在，反馈中应包含失败类型和未完成的子目标。

模拟目录存的是英文名称\code{ML reading group}，中文精确名称查询返回空。一个不检查中间结果的固定脚本会继续访问不存在的结果，或猜一个ID；带反馈的分工则将“身份未解析”退回规划器，保留用户条件，改按确认过的活动编号搜索。这次变化只重写对象查找，不必丢弃有效的目标时间与“不通知”约束。

\subsubsection{判断、行动与观察交替：ReAct}
\label{sec:nlp-execution-react}

ReAct在行动之间加入简短的当前判断，再将外部观察接入下一次选择。输入由目标、工具说明、示例和截至当前的轨迹组成；原文的提示式实现使用固定语言模型，不在每次工具返回后更新权重。可将动作选择写为$a_t\sim\pi_{\theta}(\,\cdot\mid g,h_t)$，其中$g$为目标，$h_t$包含过去的行动与真实观察。判断文字帮助组织缺失信息和下一步，但不作为环境事实存储。

算法\ref{alg:nlp-execution-react}在这种交替组织上补充本章的预算与异常契约，不将这些工程条件全归为原论文设计。终止请求也必须转换成读取或测试，写入结果未知时则优先恢复状态，不允许另一个写动作越过它。

\begin{algorithm}[带预算与结果核验的行动循环]
\label{alg:nlp-execution-react}
\begin{algorithmic}[1]
\Require 目标$g$，工具目录，初始状态$S$，调用预算$B$，决策与时间上限
\Ensure 已核验完成，或附当前证据的未完成状态
\While{决策和时间预算未耗尽}
  \State 接收用户更正，更新目标版本，标记受影响计划和事实
  \If{$B=0$}
    \State \Return 预算耗尽，附未核验项与结果未知的操作
  \EndIf
  \If{存在结果未知的写入}
    \State $a\gets$查询原请求状态或读取实际对象
  \Else
    \State 根据$g,S$形成判断并选择$a$；完成请求改为结束核验
  \EndIf
  \State 检查参数、身份、权限与前置条件；不符则有限次修正或退出
  \State $o\gets$实际调用工具$a$；$B\gets B-1$；保存参数、返回和时间
  \State 按事件更新$S$；写入超时标为未知，不能标为失败或成功
  \If{返回临时故障、状态冲突或信息不足}
    \State 按分类执行有限重试、重新读取、澄清或终止
  \EndIf
  \If{最新核验观察满足$g$及保持条件，且无未知操作}
    \State \Return 已核验完成及其证据
  \EndIf
\EndWhile
\State \Return 达到上限，附已发生变化与剩余条件
\end{algorithmic}
\end{algorithm}

本地正常轨迹确实执行了六次调用。第1次按中文名称查询返回\code{[]}；第2次保留用户\code{U7}，改用活动编号\code{ML-0418}，取得\code{R17}；第3次读取记录，观察到旧时间、旧地点和版本1；第4次查询目标时段，返回可用；第5次用代码\ref{code:nlp-execution-tool-request}更新，回执为已提交、版本2；第6次读取\code{R17}，看到4月18日14:00与海棠楼201，备注未变。独立判定器还确认\code{R18}及通知记录未变。

这些观察来自模拟服务实际返回，而“中文别名未命中，改用业务编号”是根据返回作出的控制决定。轨迹展示了ReAct式交替接口，动作由确定性程序选择，不能据此估计ReAct模型成功率。若把第1次空结果写成“已找到”，后续过程在文本上仍可能流畅，却已经失去对象证据；若调用上限为4，执行器必须停在查询之后，报告尚未修改。

\subsubsection{分层子目标与依赖图}
\label{sec:nlp-execution-subgoals}

复杂任务可以把“完成预约变更”分成解析对象、确认可用资源、写入和验收，再以有向边表示结果依赖。设$u$为目标确认，$r$为记录读取，$q$为时段查询，$w$为写入，$v$为核验，则$u\to r$、$u\to q$、$r\to w$、$q\to w$、$w\to v$给出必要顺序。只读且互不依赖的$r,q$可以并行；$w$必须等到两者都成功，不能按语言排列顺序猜测哪个先完成。

为展示多个地点，另构造会前准备预约：用户明确允许海棠楼202或203，不改变读书会在201的地点。两个地点查询节点$q_{202},q_{203}$都依赖日期确认，却不互相依赖；选择节点读取各自的真实可用结果，再把选定房间传给写入。若202不可用、203可用，顺序查询需要两次调用，并行查询也需要两次调用，只可能缩短等待时间，不能把它说成减少了调用数。若202一查即符合优先规则，顺序策略可以不查询203，而预先并行会多付一次成本。

用户更正日期后，原对象ID仍可保留，但两个时段查询和依赖它们的选择结果都应失效；若只有202查询故障，不必重做已确认的身份解析。维护这种失效传播需要为节点保存输入版本、输出证据和完成状态。分层规划的监督可来自任务与子目标轨迹，或由现有模型据工具契约生成；本例依赖图是人工构造。比较分层与直接行动时，固定初态、信息和总预算，分别统计成功率、无效调用、重规划及成本，不能从一张更完整的计划图推断收益。

\subsection{观察、状态与任务记忆}
\label{sec:nlp-execution-state-memory}

轨迹变长后，执行器需要有选择地保留信息。当前状态不是对历史文字的任意摘要，而是对任务仍有影响的事实、假设、目标和操作结果。一个字段应能回答“谁说的，针对哪个对象，在哪个版本成立”；否则历史中的正确信息也会在后续成为错误依据。存储实现与通用记忆学习超出本章范围，以下只规定执行接口。

\subsubsection{事件更新与结构化任务状态}
\label{sec:nlp-execution-state-update}

可把状态组织为$S=(g^{(v)},I,F,H,O)$：$g^{(v)}$是第$v$版用户目标，$I$为已确认对象身份，$F$为带来源和版本的观察事实，$H$为未确认假设，$O$为操作及其状态。每次事件携带来源、对象ID、观察时间、环境版本或请求号；更新函数$S_{t+1}=U(S_t,e_t)$由接口规则决定，不需要训练一个模型来猜写入是否发生。

正常轨迹在第3次调用后保存\code{R17}版本1以及旧字段，操作状态仍为“未提交”；第5次得到请求回执后保存“请求已提交，回执版本2”，但当前记录仍待回读；第6次才得到新的字段快照。写入若超时，$O$记为“结果未知”，不把期望字段复制进$F$。这种区分使后续控制知道该继续行动还是先回查。

若其他入口在第5次和第6次之间修改了备注，版本2的回执并不描述最新记录。版本3的读取应覆盖相应字段快照，同时保留版本2回执用于审计。相反，迟到的版本2响应不能覆盖已知版本3；没有可靠版本号时，需要请求顺序、服务器时间和再次读取来消除歧义，单凭本地接收时间不足以确定新旧。用户说“改到19日”更新的是目标，不意味着环境已处于19日，两者应进入不同字段。

\subsubsection{历史摘要与外部记忆检索}
\label{sec:nlp-execution-history}

完整轨迹直接输入可以保留细节，却占用上下文并引入过期信息。结构化摘要应至少保存目标版本、对象ID、否定条件、已确认字段、未知操作和原始结果定位；外部记录检索则按当前对象与子目标取得相关历史，补回尚未放入摘要的细节。关键词、BM25或嵌入负责召回，身份、时间与有效性过滤负责判断记录是否适用，检索机制沿用第\ref{chap:nlp-classification-matching-organization}章。

本例若摘要只写“预约改到18日”，就丢掉了“不发送通知”和“保留备注”；执行器可能选择一个会通知全员的更新接口。若按同名标题召回\code{R18}上次成功更新的记录，又会把另一个活动的ID带入当前任务。较可靠的摘要是“用户U7，活动ML-0418，目标18日14:00与201，不通知；R17版本1已读取；写入尚未提交”，并附各字段的来源事件号。

摘要生成可以由固定模型完成，监督若存在则来自轨迹与摘要配对；本例不训练摘要器，而用必保字段规则检查信息是否被删。遇到不确定内容，摘要应保留“待核对”，不得因压缩把“可能是R17”变成“已确认R17”。评价应同时检查关键条件保留率、所检索记录的身份正确率和执行结果；检索命中一条相似历史，只说明找到了文本，不说明现在可以重用当时的动作。

\subsubsection{反馈反思与回合记忆：Reflexion}
\label{sec:nlp-execution-reflexion}

一次失败之后，保存原轨迹未必足以让下一次尝试避开同一错误。Reflexion用评价反馈和轨迹生成一段可供后续使用的反思文字，存入回合记忆，再把记忆作为下一次行动的条件。原文区分行动者、评价器与反思器三个职责，评价可以来自规则、环境奖励或模型；它们不要求对应三个分别训练的新模型。这种适应改变的是上下文记忆，不是语言模型权重，方法依据采用NeurIPS 2023正式版本。

设第$j$次尝试的轨迹为$\tau_j$，评价为$r_j$，反思器生成$m_j=f_{\psi}(\tau_j,r_j,M_j)$，再以容量受限的更新得到$M_{j+1}$。下一次行动者条件化于原目标、新观察和$M_{j+1}$。原文区分本轮轨迹的短期记忆与跨轮反思的长期记忆，并限制保留条数；容量、选择规则和尝试次数均属于预算。评价认为通过或尝试上限到达时退出，不因反思文本写得完整便停止。

构造旧记忆含有已失效的\code{R404}：本地接口返回不存在，反思可写成“上轮直接复用了历史ID，未核对当前活动；下一轮按用户与活动编号重新解析ID，读版本后再修改”，并链接到失败调用。下一轮只有真的先搜索、得到\code{R17}并完成核验，才支持“避免了重复错误”。不保存记忆容易重试旧ID；只保存原轨迹需要行动者自己定位关键错误；反思提供压缩后的修正规则，但也可能错误归因于权限不足而建议换账号。

因此，外部错误码与反思器对原因的猜测要分别标注。模型模拟的反馈不能写成工具返回，反思提出的经验还应核对接口和实际轨迹。比较三种记忆条件时固定模型、初态和总尝试预算，并观察后续动作是否改变，不以“记住了”“吸取教训”等文字作为改进指标。

\subsection{用户协作与执行恢复}
\label{sec:nlp-execution-recovery}

执行中断不只有一种原因。参数错误应返回生成器修正字段；明确的暂时不可用可以按约定次数和退避间隔重试；版本冲突需要重新读取、检查最新条件并重规划；缺少信息或存在多个候选对象则交给用户澄清。每条恢复路径都设置退出条件，不能把所有异常都接到“再试一次”。对话中的澄清方式回用第\ref{chap:nlp-question-answering-dialogue}章，这里新增的是已发生变化如何处理。

写操作超时尤其需要单独状态。它可能尚未提交，也可能已提交但回执丢失；两者对客户端都表现为没有确定响应。控制器先记为“结果未知”，查询原请求的状态，再读取当前对象。模拟服务提供\code{request\_id}，相同请求号与相同参数的重放只返回原回执，实际写入一次；相同请求号却换参数会被拒绝。这称为\term{幂等键}（idempotency key），其作用依赖服务端持久保存请求与结果，不能靠客户端生成一个字符串自行成立。

在“写后超时”轨迹中，第5次调用已经把版本1改为2，但抛出超时；第6次查询请求，取得已提交回执；第7次读取，才确认目标字段。因此总调用7次、实际写入1次。“写前超时”轨迹的第6次请求查询为不存在；本地服务是同步执行、查询立即一致，故这一返回在本模拟契约下能排除仍在途的请求。第7次以同一请求号和同一参数重试，第8次回读，也只写入1次。真实异步系统的“暂未查到”未必表示未执行，应继续有限回查或报告未知，不能照搬本模拟的判断。

更复杂的轨迹同时包含用户更正。第5次写入18日成功但回执超时，用户随后明确把\emph{个人预约}改为19日15:00，并从另一入口把备注改成“带电脑”；活动通知本身仍是18日14:00。状态先把目标更新为第2版，却保留旧请求的未知状态。第6次查询旧请求确认已提交，第7次读取版本3，见到18日的预约和用户新备注；第8次检查19日时段，第9次以新请求号、版本3只修改时间和房间，第10次读取版本4，确认19日15:00以及“带电脑”。

这次恢复共10次调用、2次有不同目标的写入，没有重复发送通知或覆盖新备注。重用旧幂等键提交19日会产生参数冲突；盲目重发旧请求则可能恢复已经过时的18日目标；只把回复改成“已改到19日”不会改变实际记录。支持条件写入的版本检查可以阻止覆盖并发更新，但冲突后仍要重新绑定参数，而不是把旧快照整条写回。

$\tau^2$-Bench以双方都能影响环境的设置强调了用户行动对协作任务的作用，相关来源见文献说明。本例只借用“共同状态会变”的问题，不复现该基准的环境或成绩。任务如果要求撤销已发生的变化，应核对是否存在明确逆操作及其副作用；本例可以在最新版本上重新设置个人预约字段，却不能据此推断付款、发信或抢占资源都能撤销。没有可靠逆操作时，应据已发生状态重新求解并告知用户，而不是假装回到初态。

\subsection{完成判断与执行评价}
\label{sec:nlp-execution-completion}

结束判定器应独立读取环境，把“已完成”翻译成可检查条件。本例要求正确对象的目标时间与地点成立、备注遵循最新用户更正、另一预约不变、未发送通知。若$g_j(s)$为第$j$项目标或保持条件，终态判定可写为
\[
 J(s)=\prod_{j=1}^{m}\mathbb{I}[g_j(s)].
\]
各条件从任务规格和受控读取取得，不由待评价执行器根据自己的回复生成。乘积要求所有条件共同成立；只检查时间就会放过地点仍旧的部分完成。若过程另有限制，例如必须先获授权再修改，则单独检查事件顺序，不能从终态反推整个过程合规。

WebArena对信息查询比较答案，对导航、内容与配置任务使用定位器读取相关页面或存储状态，再检查所需内容；它并不要求行动序列与某条参考轨迹逐字一致。其方法也提醒我们区分答案判定与环境判定：两种不同操作路径可以达到相同合格结果，相同按钮点击却可能作用于不同对象。本例用直接读取模拟记录及审计记录代替网页定位器，不把它当作真实WebArena实验。

表\ref{tab:nlp-execution-outcomes}给出本地保存的完整运行记录的摘要。每条记录含初态、参数、真实返回、最终记录及独立判定；成功场景只有1次业务写入，查询与核验也照常计费。“预算5”场景尤其值得区分：环境已被正确修改，但执行器没有预算做第6次回读，只能报告“已提交，尚未核验”，独立评价器随后能看到目标确实达到。环境结果与执行器当时有权声称的结果不是同一个字段。

\begin{table}[htbp]
  \centering
  \begin{tabular}{lccc}
    \toprule
    轨迹 & 调用数 & 写入数 & 外部目标判定 \\
    \midrule
    正常执行 & 6 & 1 & 通过 \\
    只改时间、未改地点 & 1 & 1 & 不通过 \\
    读取服务不可用 & 3 & 0 & 不通过 \\
    预算为4，查询后停止 & 4 & 0 & 不通过 \\
    预算为5，未做回读 & 5 & 1 & 通过，未自验 \\
    成功修改同名另一对象 & 1 & 1 & 不通过 \\
    只宣布完成、未调用 & 0 & 0 & 不通过 \\
    \bottomrule
  \end{tabular}
  \caption{模拟执行的结果与成本。所有输入和轨迹均为教学构造；独立评价器读取不计入执行器调用预算。目标判定不隐藏环境故障或部分完成，停止原因另保存在记录中。}
  \label{tab:nlp-execution-outcomes}
\end{table}
\FloatBarrier

比较真实系统时，固定环境版本、初态、资料、工具权限、模型和预算，再分别统计工具选择、参数绑定、状态更新和终态成功。随机系统需重置环境后重复运行，报告成功比例及不确定性，同时报告调用、时延、重复副作用和需要用户澄清的次数。参数错误比例高，应先检查对象与字段绑定；参数正确但终态失败，则继续检查并发变化、执行异常和完成判定。不能只把所有失败平均成一个分数后猜原因。

多次环境尝试至少一次完成的概率、每次完整运行的成功比例，以及选择器最终交付一次的完成率具有不同分母；它们也不自动等同于程序生成的pass@$k$。若环境不能无损重置，第二次尝试已经受第一次副作用影响，必须记录这种依赖，不能仍按独立同分布试验估计。评价不仅需要问“最后成功了吗”，还要说明它经历了什么变化、消耗了多少资源，以及哪些事实在停止时仍然未知。

\section*{文献说明}
\label{sec:nlp-execution-literature}

以下按方法用途列出来源，便于将正文中的机制与原始研究对照。数量题、约束树、程序补丁和预约轨迹均为教学构造；它们用于解释和核验计算，不复现论文的实验结果。

\begin{fullwidth}
CoT的分步提示、self-consistency的路径采样与答案聚合，分别见以下研究。
\par\fullcite{nlp-execution-wei2022cot}
\par\fullcite{nlp-execution-wang2023consistency}
\end{fullwidth}

\begin{fullwidth}
ToT给出状态扩展、评价与BFS／DFS控制；PAL给出程序化中间步骤与解释器执行。正文的固定搜索树、操作计数和金额程序另行构造。
\par\fullcite{nlp-intro-yao2023tot}
\par\fullcite{nlp-execution-gao2023pal}
\end{fullwidth}

\begin{fullwidth}
GSM8K验证器支持结果监督与候选重排的讲解；过程监督论文支持步骤标签及正标签概率乘积。正文手设分数不取自其模型输出。
\par\fullcite{nlp-intro-cobbe2021}
\par\fullcite{lightman2023verify}
\end{fullwidth}

\begin{fullwidth}
Z3命令依据官方指南，Lean命题与证明项、公理检查依据官方教材；Python AST接口依据与教学运行对应的3.9版本文档。
\par\fullcite{nlp-execution-z3guide}
\par\fullcite{nlp-execution-lean4}
\par\fullcite{nlp-execution-python-ast}
\end{fullwidth}

\begin{fullwidth}
HumanEval的测试协议与组合估计量来自下列工作。它支持pass@$k$的定义，不为自适应补丁修订或实际选择器的单次成功率提供等价关系。
\par\fullcite{nlp-intro-chen2021}
\end{fullwidth}

\begin{fullwidth}
\raggedright
SynCode的解析状态、DFA掩码及保证方向按第2版正文说明；Self-Debugging的候选解释与执行反馈按所引版本说明。
\par\fullcite{nlp-execution-ugare2024syncode}
\par\fullcite{nlp-execution-chen2023selfdebug}
\end{fullwidth}

\begin{fullwidth}
SWE-bench说明基于真实仓库问题的补丁评价；Agentless第2版支持分层定位、局部修复和补丁验证三阶段。两文件日期仓库与严格的教学选择器并非其原始实验。
\par\fullcite{nlp-intro-jimenez2024}
\par\fullcite{nlp-execution-xia2024agentless}
\end{fullwidth}

\begin{fullwidth}
ReAct支持判断、行动与观察交替；Reflexion采用NeurIPS 2023正式版，支持反馈反思与跨回合记忆。幂等回查、版本控制和本章异常协议是配套模拟的接口设计。
\par\fullcite{yao2023react}
\par\fullcite{nlp-execution-shinn2023reflexion}
\end{fullwidth}

\begin{fullwidth}
$\tau^2$-Bench支持用户与执行器共同改变环境的协作视角；WebArena支持按答案或环境功能结果进行评价，而非只匹配参考动作序列。
\par\fullcite{nlp-intro-barres2025}
\par\fullcite{nlp-intro-zhou2024webarena}
\end{fullwidth}

\section{本章小结}
\label{sec:nlp-execution-summary}

语言要求能否落实，取决于是否先确定可交付对象，再选择与之对应的检查方式。金额题要把数量、单位和折扣作用域恢复成正确表达式；约束题要区分找到可行解、搜索未找到与已证明无解；程序任务要对照规格运行候选，并在修订目标行为时保持相关旧行为。环境任务又多了一项责任：记录已发生的副作用，处理变化的用户要求和共同状态，最后读取实际结果。

候选形成与反馈机制可以组合，但承担的职责不同。CoT提供可检查步骤，自一致性聚合答案，ToT分配状态搜索预算，PAL把计算交给解释器；ORM和PRM帮助选择候选，规则与形式检查则在各自前提下直接验证特定对象。程序修订把测试失败接回修改，ReAct把外部观察接回下一行动，Reflexion把跨轮经验接回上下文。这些组合的收益都需要与额外候选、调用和核验成本一起衡量。

章首“预约是否真的改好”的问题，不能由一句完成声明回答。应当能指出改了哪条记录、使用了哪版目标、保留了哪些字段、读到了什么结果；写入超时而尚未回查时，正确交付可以是清楚说明未知与后续核验。下一章转向语音任务，输入输出增加了声学信息和时间约束，但取得证据、保存过程与验收目标这一区分仍然适用。
